% LiteraLevel: guion actualizado a Angular + API, 14 septiembre 2026.
% Fuente hablada equivalente: guion_video.md; contador: tools/contar_guion.py.
\documentclass[11pt,a4paper]{article}
\usepackage{fontspec}
\usepackage[spanish,es-noquoting,es-nodecimaldot]{babel}
\usepackage[a4paper,margin=2.4cm,headheight=14pt]{geometry}
\usepackage{graphicx}
\usepackage{microtype}
\usepackage{setspace}
\usepackage{fancyhdr}
\usepackage{titlesec}
\usepackage{enumitem}
\usepackage{booktabs}
\usepackage{tabularx}
\usepackage{xcolor}
\usepackage[most]{tcolorbox}
\usepackage[hidelinks]{hyperref}
\definecolor{llazul}{HTML}{1F4E79}
\definecolor{llcian}{HTML}{2E86AB}
\definecolor{llgris}{HTML}{5A5A5A}
\definecolor{llgrisclaro}{HTML}{E8ECF0}
\definecolor{llfondo}{HTML}{F7F9FB}
\setstretch{1.10}
\setlist{itemsep=4pt,topsep=4pt}
\pagestyle{fancy}
\fancyhf{}
\lhead{\small\color{llgris} LiteraLevel · guion del vídeo}
\rhead{\small\color{llgris} Módulo 5 · 20 pts}
\cfoot{\small\thepage}
\titleformat{\section}{\large\bfseries\color{llazul}}{}{0em}{}
\newtcolorbox{dice}{colback=white,colframe=llazul!55,boxrule=0.8pt,arc=1mm,
  left=8pt,right=8pt,top=7pt,bottom=7pt,enhanced,breakable}
\newtcolorbox{ve}{colback=llfondo,colframe=llgrisclaro,boxrule=0.7pt,arc=1mm,
  left=8pt,right=8pt,top=6pt,bottom=6pt,enhanced,breakable}
\newcommand{\fase}[2]{\section{#1}\noindent{\small\bfseries\color{llcian}#2}\par\medskip}
\begin{document}
\hypersetup{pageanchor=false}
\begin{titlepage}
\centering
\begin{minipage}{0.22\textwidth}\raggedright
\includegraphics[height=2.2cm]{logos/logo_unam.png}
\end{minipage}\hfill
\begin{minipage}{0.52\textwidth}\centering
{\large\bfseries Universidad Nacional Autónoma de México}\par\medskip
{\normalsize Facultad de Estudios Superiores Acatlán}
\end{minipage}\hfill
\begin{minipage}{0.22\textwidth}\raggedleft
\includegraphics[height=2.2cm]{logos/logo_fes_acatlan.png}
\end{minipage}
\vspace{2cm}
{\Huge\bfseries\color{llazul} Guion del vídeo}\par\bigskip
{\LARGE LiteraLevel}\par\bigskip
{\large Angular + API FastAPI}\par
\vspace{1.5cm}
{\Large Miguel Reyes Carrillo}\par\bigskip
{\large Proyecto integrador · Módulo 5}\par
{\normalsize Diplomado en Ciencia de Datos · Generación 34}\par\bigskip
{\normalsize Mtro. Fernando Barranco Rodríguez}\par
\vfill
\begin{ve}
\centering
\textbf{MP4 · máximo 5:00 · cara visible todo el tiempo}\par\medskip
Objetivo: \textbf{4:40}. Texto hablado: \textbf{480 palabras}.\par
3:12 a 150 palabras/min; unos 3:42 a 130.\par
Con hasta 55 segundos de acciones: presupuesto conservador \textbf{4:37}.\par\medskip
\small Estimación, no ensayo cronometrado. Actualizado el 14 de septiembre de 2026.
\end{ve}
\end{titlepage}
\hypersetup{pageanchor=true}

\section{Preparación de la toma}
El producto que se muestra es \textbf{Angular con la API FastAPI}. Gradio
queda como motor y respaldo. OBS debe mostrar la app y la cara de Miguel
durante toda la grabación, sin superponer la cámara a los resultados.
\begin{itemize}
\item Abrir \url{https://mikemaki1763-literalevel-front.static.hf.space/} a
390 px de ancho o duplicar el móvil. Comprobar que texto y cámara son legibles.
\item Entrar en Ajustes antes de grabar. No mostrar contraseña, identificador
de lector, consola, secretos ni paneles de administración.
\item Preparar un perfil \textbf{ilustrativo} de 1,20 sobre 2, buscar
\emph{dostoyevski} y abrir \emph{Crime and Punishment} (PG2554). No es una
sesión real registrada ni un resultado fijo del cuestionario.
\item Abrir previamente la curva guardada. El ejemplo documentado ronda
1,70; si la medición vigente cambia, decir el valor visible. No mezclar el
ensayo de EPUB (203.699 palabras, 840 ventanas) con el de TXT (203.415, 836),
ni prometer esos tiempos en el móvil.
\item Preparar una palabra y una frase corta. \textbf{Traducción y explicación
llegan con la etapa de IA}, no con una supuesta respuesta completa en 40 ms.
Conservar los avisos de acepción dudosa. Sin permiso, mostrar la parte gratuita.
\item Ensayar micrófono y audio de la app. No prometer un fonema concreto:
la corrección depende de lo que se entienda. Anunciar un error intencional si
se representa como demostración.
\item Comprobar el apoyo articulatorio en la ficha IPA y en la corrección:
el modelo Blender de 38 piezas se carga al abrir el 3D y conserva un respaldo
2D. Es una ayuda didáctica, no clínica. Ensayar su apertura sin alargar la toma.
\item Las portadas FLUX siguen pendientes. No fingir imágenes generadas,
respuestas en vivo ni tiempos de espera que no ocurrieron.
\item Preparar un respaldo del mismo recorrido. Si se usa, rotular
\textbf{grabación previa}; no presentarlo como respuesta en vivo.
\end{itemize}

\section{Ritmo y ensayo}
ANSVA organiza el relato: atención, necesidad, satisfacción, visualización y
acción. Reservar hasta 55 segundos para clics, frase leída, audio y pausas.
Si el ensayo supera 4:40, recortar acciones; no acelerar hasta resultar
incomprensible. La presentación es apoyo para la defensa: no hay que leer
todas las diapositivas en el vídeo.

El contador oficial es \texttt{tools/contar\_guion.py}. Cuenta los parlamentos
de \texttt{guion\_video.md}, sin notas de producción. El ensayo completo con
acciones es lo que decide si cabe en cinco minutos.
\newpage

\fase{A · Atención}{0:00 - 0:25}
\begin{ve}\textbf{En pantalla:} LiteraLevel. Miguel visible. Pausa breve tras
la pregunta.\end{ve}
\begin{dice}
¿Te ha pasado? Abres un libro en inglés que te interesa y, a las pocas
páginas, lo dejas. Quizá no era falta de capacidad: quizá necesitabas otra
lectura. LiteraLevel te ayuda a comparar la dificultad de un libro con tu
nivel antes de empezar, y te acompaña mientras lees.
\end{dice}

\fase{N · Necesidad}{0:25 - 0:50}
\begin{ve}\textbf{En pantalla:} ficha de libro con nivel y curva por
capítulos.\end{ve}
\begin{dice}
Hay miles de libros gratuitos, pero elegir uno adecuado sigue siendo difícil.
Además, un libro no tiene una sola dificultad: sus capítulos cambian.
Nuestro usuario es un hispanohablante que quiere leer en inglés con autonomía.
Necesita una orientación comprensible, no otro examen ni una etiqueta que
prometa más de lo que puede medir.
\end{dice}

\fase{S · Satisfacción}{0:50 - 3:35}
\begin{ve}\textbf{En pantalla:} Angular. Recorrer Tu nivel, Encontrar, ficha,
lector, Pronunciar y Mis palabras. Mostrar el perfil ilustrativo preparado;
las credenciales ya están introducidas fuera de toma.\end{ve}
\begin{dice}
Esta es la aplicación. La interfaz está hecha en Angular y se comunica con
una API FastAPI. El motor de Python también conserva Gradio como respaldo.

Primero, sitúo mi nivel con un cuestionario o con un texto que ya leo cómodo.
Para este ejemplo preparé un perfil de \textbf{1,20 sobre 2}, banda intermedia.

Ahora busco. Puedo describir en español lo que me apetece leer. Escribo
\emph{dostoyevski} y elijo \emph{Crimen y castigo}. La ficha muestra datos del
texto y su curva de dificultad. En este ejemplo, el libro ronda \textbf{1,70}.
Al compararlo con mi perfil, veo que me exigirá más.

Entro al lector y toco una palabra. Primero aparece el diccionario, con la
pronunciación disponible. Después, si tengo permiso para usar IA, llegan la
traducción y la explicación de esa frase. El diccionario y el texto generado
se distinguen; si la acepción es dudosa, la aplicación lo avisa.

También puedo practicar una frase en voz alta. Grabo y pido la corrección.
Comparo las palabras que se entendieron con las que quería decir, escucho el
audio de referencia y abro el apoyo de articulación para practicar.

Las palabras consultadas quedan en \textbf{Mis palabras}, con su frase, para
volver a ellas sin perder el contexto.

Detrás está un modelo distilroberta afinado en regresión. Su validación interna
registró un Spearman de \textbf{0,908}. Es una medida de orden, no una nota para
el lector. \textbf{No certificamos ningún nivel del MCER}: usamos una escala
propia de cero a dos para orientar la lectura.
\end{dice}

\fase{V · Visualización}{3:35 - 4:15}
\begin{ve}\textbf{En pantalla:} 53.179 libros de Gutenberg y 585.111 capítulos
medidos. No confundir con el catálogo objetivo de 61.642 ni con los corpus
graduados.\end{ve}
\begin{dice}
El catálogo medido reúne \textbf{53.179 libros} de Project Gutenberg y
\textbf{585.111 capítulos}. La aplicación aprovecha esa información para
comparar lecturas y, cuando hay capítulos adecuados, sugerir por dónde empezar.

Leer, buscar y consultar el diccionario no requieren llamadas de pago a
OpenAI. La explicación y la voz tienen su puerta de acceso y su cuota.
La siguiente validación será comprobar con lectores si esta ayuda mejora
su experiencia y continuidad.
\end{dice}

\fase{A · Acción}{4:15 - 4:40}
\begin{ve}\textbf{En pantalla:}
\url{https://mikemaki1763-literalevel-front.static.hf.space/}.
Mantener la mirada un segundo al terminar; cara visible hasta el final.\end{ve}
\begin{dice}
Este es el enlace de LiteraLevel. Se abre en el navegador del móvil o del
ordenador. Elige un libro que te interese, compara su dificultad con tu perfil
y prueba una página con las ayudas de lectura. Ese es el siguiente paso:
no demostrar cuánto inglés sabes, sino encontrar qué te apetece seguir leyendo.
\end{dice}
\newpage

\section{Comprobación antes de grabar}
\noindent\begin{tabularx}{\textwidth}{@{}p{0.22\textwidth} X@{}}
\toprule
\textbf{Comprobación} & \textbf{Criterio} \\
\midrule
Duración real & Ensayo completo menor de 4:40; MP4 final nunca supera 5:00 \\[5pt]
Cámara & Cara visible y reconocible todo el tiempo \\[5pt]
Pantalla & Texto legible; cámara no tapa controles ni resultados \\[5pt]
Privacidad & Sin contraseñas, tokens, IDs de lector ni datos de terceros \\[5pt]
Sonido & Micrófono y audio de la app se oyen sin saturar \\[5pt]
Estado & Front abre y backend responde; no basta un 200 con HTML \\[5pt]
Ejemplo & Libro y palabra preparados; valores dichos coinciden con pantalla \\[5pt]
IA & No atribuir traducción al diccionario ni ocultar el aviso de duda \\[5pt]
Respaldo & Si se reproduce una captura, rotular grabación previa \\[5pt]
Formato & Exportar MP4 y reproducirlo completo antes de entregar \\
\bottomrule
\end{tabularx}
\bigskip

Las métricas corresponden a la documentación del modelo, no a una nueva
evaluación hecha al grabar. Este guion no demuestra mejora pedagógica ni
ausencia total de fallos; esos resultados necesitan una evaluación propia.
\end{document}
